\section{Conectividad de las redes}
\label{sec:discusion-conectividad-v2}

Los dos indicadores de la Sección~\ref{sec:conectividad} —clases de salida desconectadas y conexiones inútiles— responden a mecanismos distintos y conviene discutirlos por separado, aunque ambos comparten una misma raíz: ninguno de los operadores de dispersión de MOEACKF ni de S-NSGA-II tiene noción alguna de a qué neurona o a qué salida pertenece cada sinapsis, y ninguno de los dos objetivos de la optimización ($f_1$, complejidad; $f_2$, error de entrenamiento) penaliza directamente una sinapsis activa que no forma parte de ningún camino completo entre la entrada y la salida.

\subsection{Clases de salida desconectadas: desbalance de clases}

La asimetría por clase ya señalada en el análisis de resultados —en \texttt{ds2} es siempre la clase 0 la que pierde su conexión, en \texttt{ds3} es siempre la clase 1— coincide exactamente con la clase minoritaria de cada dataset: en \texttt{ds2} (Climate) la clase 0 representa apenas 46 de 540 instancias (8.5\%); en \texttt{ds3} (German) la clase minoritaria representa 300 de 1000 (30\%). En \texttt{ds1} (Australian, clase 0/clase 1: 55.5\%/44.5\%) y \texttt{ds4} (Sonar, 53.4\%/46.6\%), donde ambas clases están razonablemente balanceadas, el porcentaje de redes con una salida desconectada es nulo o casi nulo (0\% en tres de las cuatro combinaciones, 6.5\% en la restante). Esto es consistente con que, al no distinguir los operadores de dispersión entre clases al apagar pesos, basta con que $f_2$ (el error de entrenamiento) recompense predecir siempre la clase mayoritaria —una estrategia que en \texttt{ds2} ya alcanza 91.5\% de exactitud sin necesitar ninguna conexión hacia la clase minoritaria— para que la búsqueda converja hacia soluciones que sacrifican por completo la conectividad de esa clase, sin que el objetivo de exactitud lo impida.

La excepción parcial es \texttt{ds3}$\times$SNSGAII, con 0\% de salidas desconectadas pese al mismo desbalance 70/30 que en \texttt{ds3}$\times$MOEACKF produce 80.6\%. La diferencia es coherente con el diseño de cada operador: el operador de mutación de dispersión de S-NSGA-II muta un nivel de dispersión objetivo por individuo y apaga o reactiva pesos al azar hasta alcanzarlo, mientras que MOEACKF construye su máscara inicial activando un número aleatorio $k \in [1, D]$ de variables mediante selección por torneo sesgada por aptitud, sin ninguna garantía de que ese conjunto —que puede ser tan pequeño como una sola sinapsis de las $D$ totales— toque la neurona de salida de la clase minoritaria. Con un espacio de decisión de cientos o miles de variables, es más fácil que un sorteo casi puntual dependa del azar de la clase minoritaria que un ajuste incremental de densidad la deje completamente fuera, aunque esta lectura es una asociación entre el mecanismo del operador y el patrón observado, no una prueba directa sobre el código en ejecución.

\subsection{Conexiones inútiles: asimetría estructural y presión de poda débil}

El porcentaje de conexiones inútiles no aparece asociado con el desbalance de clases sino con la dimensionalidad de entrada: promediando ambos algoritmos por dataset (y excluyendo \texttt{ds2}$\times$SNSGAII por la razón que se explica más abajo), el porcentaje sube de 31.5\% en \texttt{ds1} (14 características) a 54.3\% en \texttt{ds3} (24 características) y 57.3\% en \texttt{ds4} (60 características). Esta tendencia es consistente con una explicación estructural propia de la codificación: cada neurona oculta tiene $n_F+1$ posiciones de entrada posibles (una por característica más el sesgo) pero solo $n_O=2$ posiciones de salida posibles, de modo que, ante cualquier mecanismo de activación de pesos sin conciencia de la estructura del grafo, la probabilidad de que una neurona reciba al menos una entrada activa crece con $n_F$ mucho más rápido que la probabilidad de que además tenga una salida activa —fija en solo dos posiciones—. Cuantas más características tiene el dataset, más fácil es entonces que una neurona termine con entrada real pero sin salida (sumidero) que con ambas (viva), lo que aumenta la fracción de sinapsis activas atrapadas en caminos incompletos. Esta lectura es razonable con solo cuatro datasets y ocho combinaciones, y no se contrastó de forma independiente del nivel de dispersión de cada solución.

\texttt{ds2}$\times$SNSGAII es la excepción que confirma esa lectura por la vía opuesta: su solución de mejor exactitud alcanza 92.1\% de accuracy con solo el 0.32\% de las 840 sinapsis activas y una única neurona oculta viva. Con tan pocas sinapsis activas en total, no queda margen combinatorio para que aparezcan neuronas parcialmente conectadas, y el 0\% de conexiones inútiles es entonces un efecto secundario de la extrema dispersión y no evidencia de una red mejor estructurada que el resto.

A esta asimetría estructural se suma una segunda causa, de naturaleza algorítmica: ninguno de los dos objetivos premia explícitamente eliminar una conexión inútil. Retirar una sinapsis que llega a una neurona sumidero (o que sale de una neurona fuente) reduce $f_1$ sin alterar $f_2$ en absoluto, de modo que la solución resultante domina estrictamente a la original y en teoría debería reemplazarla en el frente. En la práctica, esa poda solo ocurre si algún operador de variación llega a generar exactamente esa vecina —con esa sinapsis puntual desactivada y ninguna otra— dentro del presupuesto de evaluaciones disponible, algo cada vez menos probable a medida que $D$ crece (de 680 en \texttt{ds1} a 2520 en \texttt{ds4}) y que ninguno de los operadores dirige la búsqueda hacia esa vecina en particular. La diferencia entre algoritmos observada en la Tabla~\ref{tab:conectividad-algo} (49.5\% en MOEACKF frente a 32.7\% en S-NSGA-II) es consistente con esta lectura: la máscara de MOEACKF activa posiciones sueltas del vector de decisión sin agrupación alguna, mientras que la máscara de S-NSGA-II opera en franjas contiguas que —dado que la codificación agrupa las entradas de cada neurona oculta en un bloque contiguo del vector— tienden a activar o desactivar el bloque completo de una neurona en conjunto, produciendo con algo más de frecuencia neuronas enteramente vivas o enteramente muertas en vez de parcialmente conectadas.

\subsection{Implicancia práctica}

Ambos fenómenos afectan directamente a la red que se desplegaría en la práctica: en la Tabla~\ref{tab:conectividad-algo}, poco más de un tercio de las soluciones de mejor exactitud de MOEACKF y poco más de un cuarto de las de S-NSGA-II son incapaces de predecir una de las dos clases sin importar la entrada, y en promedio entre un tercio y la mitad de las sinapsis activas de esas soluciones no contribuye a ninguna predicción. Ninguno de los dos indicadores se corrige evaluando de mejor manera la red resultante, sino interviniendo en los operadores que deciden qué sinapsis permanecen activas: una inicialización o un operador de reparación que garanticen al menos un camino activo hacia cada clase de salida atacarían directamente el primer problema, y una poda estructural —que retire sinapsis pertenecientes a neuronas sumidero o fuente sin costo en $f_2$— atacaría el segundo; ninguna de las dos intervenciones fue implementada ni evaluada en el marco de esta tesis.
